\documentclass[10pt,a4paper]{amsart}
\usepackage[margin=19mm]{geometry}
\usepackage{amsmath,amssymb,mathtools}
\usepackage[scr=rsfs,cal=euler]{mathalfa}
\usepackage{xfrac}
\usepackage[unicode,hidelinks,bookmarks=false]{hyperref}
\newcommand{\ie}{i.e.\ }
\newcommand{\cf}{cf.\ }
\newcommand{\resp}{respectively\ }
\newcommand{\Subsection}{\subsection}
\providecommand{\nobreakdash}{\nobreak\hbox{-}}
\setlength{\emergencystretch}{2em}
\multlinegap0pt
\usepackage{mdframed}
\usepackage{mdframed}
\usepackage{mdframed}
\usepackage{mdframed}
\usepackage{bm}
\usepackage{dsfont}
\usepackage[shortlabels]{enumitem}
\usepackage{amssymb}
\usepackage{mdframed}
\newtheorem{proposition}{Proposition}
\newtheorem{thm}{Theorem}
\usepackage{cleveref}
\crefname{proposition}{Proposition}{Propositions}
\crefname{thm}{Theorem}{Theorems}
\newcommand{\Acc}{\mathrm{Acc}_p(h)}
\newcommand{\ECE}{\mathrm{ECE}(h)}
\title{On Uncertainty Calibration for Equivariant Functions}
\author{Edward Berman, Jacob Ginesin, Marco Pacini, Robin Walters}
\date{Source: arXiv:2510.21691v4 (CC BY 4.0)}
\begin{document}
\maketitle

\noindent\emph{Source and licence.} On Uncertainty Calibration for Equivariant Functions. Version arXiv:2510.21691v4 (CC BY 4.0), distributed by arXiv under the Creative Commons Attribution 4.0 International licence, http://creativecommons.org/licenses/by/4.0/, as recorded in the arXiv licence metadata for this exact version and shown on the abstract page https://arxiv.org/abs/2510.21691v4. Full licence notice: \texttt{licenses/arxiv-2510.21691-CC-BY-4.0.txt}.\par\smallskip

\noindent\textit{Editorial scope.} Counted unit: the article's proposition prop:bilipschitzbounds (source0.tex lines 1044--1051). The article's complete original proof of it is delivered with it (source0.tex lines 1053--1070, one \texttt{proof} environment of 18 lines). No mathematics is written, repaired or re-derived: the statement and the proof are the article's own lines, in the article's own notation, and no retained line is substituted or re-flowed.  Retained uncounted as the source-local context the proof needs: lines 924--926; lines 941--945; lines 1023--1025.  Nothing was omitted from any retained span, so the package carries no omission record.  No retained line contains a citation, so the wrapper carries no bibliography and no bibliography entry.  No retained line contains a figure environment, an image-inclusion command or drawing code, so no image is delivered and the package declares no asset dependency.  Editorial formatting only, in addition: an AMS \texttt{amsart} preamble that restates the article's own declarations and macros used by the retained lines, namely the environments \texttt{proposition}, \texttt{thm} on separate counters, together with the standard packages those lines need.  The retained environments print in this document as Proposition 1, Theorem 1 in order of appearance, whereas the article prints the same items with the numbering of its own full document.  The complete native source of the article as distributed in the arXiv e-print for this exact version is bundled unchanged as \texttt{sources/187743/source0.tex}.
\begin{proposition} \label{prop:naive}
ECE is bounded from above by $\frac{1}{2} + \int_0^{1}r(p)|\frac{1}{2} - p|dp$.
\end{proposition}

\begin{thm} \label{thm:inc_bound}  Denote the fiber $\mathcal{F}_p = h_P^{-1}(p)$. Denote the total density on a fiber $\mathcal{F}_p$ by $q(\mathcal{F}_p) = \int_{\mathcal{F}_p}q(x)dx$ and the renormalized density by $q_p(x) = q(x)/q(\mathcal{F}_p)$. Let $k_p(Gx)$ be the total dissent of an orbit on $\mathcal{F}_p$ with the renormalized probability $q_p(x)$. Let $F_p$ be a fundamental domain for the action of $G$ on $\mathcal{F}_p$. Assume that $h$ is incorrectly invariant under $G$ on each fiber of $\mathcal{F}_p$. In other words, $h$ satisfies $h(gx) = h(x)$ for $g \in G$, $x \in \mathcal{F}_p$, but $f(x) \neq f(gx)$ for some $x$, $g$. Let  $P_1 = \{ p \colon \Acc \leq 1/2\}$ and $P_2 = \{ p \colon \Acc \geq 1/2\}$. Let $Gx^*$ be the orbit with the smallest nonzero total dissent $k(Gx^*)$, i.e., $x^* = \underset{x \in X}{\arg \min} ~k(Gx) $.  Then ECE is bounded above by 
\[
\ECE  \leq \frac{1}{2} + \int_0^{1}r(p)\left|\frac{1}{2} - p\right|dp -  k(Gx^*)\int_{P_2}r(p) dp. 
\] 
\end{thm}

\begin{align}
  r(p) = \int_Xq(x)\delta(p - h_P(x))dx &= \int_{\mathcal{F}_p}\frac{1}{|\nabla h_P(x)|}q(x)dx_p \label{eq:reparam_density}
\end{align}

\begin{proposition}\label{prop:bilipschitzbounds}
Assume $h_P(x)$ is differentiable, has gradient nowhere $0$, and is $(K_1, K_2)$ bi-Lipschitz continuous. Let $Gx^\diamond$ be the orbit with the least integrated probability density, i.e., $x^\diamond = \underset{x \in X}{\arg \min} \int_{Gx}q(x)dx$. 
Then
\[
\ECE \leq \frac{1}{2} + \frac{K_2}{4} + \min \left\{ 0, - k(Gx^*)K_2\int_{P_2}dp\int_{Gx^\diamond}q(x)dx \right\}.
\]
% The upper bound on ECE in \cref{thm:inc_bound} is then $$\frac{1}{2} + \frac{K_2}{4}  - k(Gx^*)K_2\int_{P_2}dp\int_{Gx^\diamond}q(x)dx$$ and the upper bound on ECE in \cref{prop:naive} is $$\frac{1}{2} + \frac{K_2}{4}.$$
\end{proposition}

\begin{proof}
We start with the bound from \cref{thm:inc_bound}. Substituting the expression for $r(p)$ from  \cref{eq:reparam_density} into the upper bound from \cref{thm:inc_bound} and using the Lipschitz constant to bound the gradient gives $$\ECE \leq \frac{1}{2} + \int_0^{1}\int_{\mathcal{F}_p}K_2q(x)dx_p\left|\frac{1}{2} - p\right|dp  - k(Gx^*)\int_{P_2}\int_{\mathcal{F}_p}K_2q(x)dx_pdp.$$ We now relate this inequality to integrals over $X$ and $Gx^\diamond$ in order to remove the dependence on $h_p$. Notice,
%
\begin{equation*}
\int_0^{1}\int_{\mathcal{F}_p}K_2q(x)dx_p\left|\frac{1}{2} - p\right|dp \leq  \int_0^{1}\int_{X}K_2q(x)dx\left|\frac{1}{2} - p\right|dp =  \frac{K_2}{4}.
\end{equation*}
%
and
\begin{equation*}
 % k(Gx^*)
 \int_{P_2}\int_{\mathcal{F}_p}K_2q(x)dx_pdp \geq 
 % k(Gx^*)
  \int_{P_2}\int_{Gx^\diamond} K_2q(x)dxdp =
 %k(Gx^*)
  K_2\int_{P_2}dp\int_{Gx^\diamond}q(x)dx. 
\end{equation*}
Thus, $$\ECE \leq \frac{1}{2} + \frac{K_2}{4} - k(Gx^*)K_2\int_{P_2}dp\int_{Gx^\diamond}q(x)dx.$$ The upper bound for \cref{prop:naive} can be derived in the same way or seen as the special case where $k(Gx^*) = 0$ or $P_2 = \emptyset$ and the improvement from \cref{thm:inc_bound} is vacuous.
\end{proof}

\end{document}
